\documentclass[12pt]{article}
\usepackage[margin=1in]{geometry}
\usepackage{enumitem}
\usepackage{titlesec}
\usepackage{parskip}
\usepackage{amsmath}
\usepackage{graphicx}
\usepackage{hyperref}
\usepackage{fontenc}

\title{\textbf{Kalyvas 2006 RG Notes}\\
\large Kalyvas, Stathis N. \textit{The Logic of Violence in Civil War}.\\
New York: Cambridge University Press, 2006.}
\author{}
\date{}

\begin{document}

\maketitle

\section{Main Idea}



%--------------------------------------------------
\section{General Notes}
%--------------------------------------------------

\begin{itemize}[leftmargin=*, itemsep=4pt]
    \item Violence in civil war is a joint product of two series of interactions: those between actors at local versus central levels, and those between combatants and noncombatants
    \item Civil war presents an opportunity for local groups to pin their own policy or ideological goals onto the movement and in particular to use denunciation to impact otherwise extraneous concerns
    \item Violence can be indiscriminate or selective, but selective violence requires information that is costly/difficult to organize at the central unit and instead must be gathered through local networks (especially those of denunciation). As a result, civil war requires some form of alliance or cooperation between centralized actors who need information and local actors who need coercive forces. (This is the driving mechanism of what differentiates  selective versus indiscriminate violence -- the willingness of the state to endorse local yet extraneous conflicts and provide protection to denouncers encourages such denunciation, which allows for selective violence.)
    \item Violence is driven not by ethnicity / ideology but by the ability to achieve control over an area and the control already in place. Indiscriminate violence appears where control is low (defectors cannot be individually identified) while selective violence appears where control is moderate (defectors can be identified but are not dissuaded from defecting in the first place)
    \item Indiscriminate violence negates any incentive for individuals to join or support the perpetrators and is subsequently counterproductive: if you cannot guarantee the safety (relatively) of an actor who supports, then you cannot convince them to support if they otherwise would not
    \item Five zones of territorial control (reflected over a center point):
    \begin{itemize}
        \item 1 \& 5: total control by a single party
        \item 2 \& 4: partial (but dominant) control by one party
        \item 3: parity and contestation; expectation of maximal violence, particularly indiscriminate
    \end{itemize}
\end{itemize}


\section{Important Specific Factual Claims}

\begin{itemize}[leftmargin=*, itemsep=4pt]
    \item Unfolding identities and alliances through civil wars are often endogenous and not necessarily reflective of previous violence or cleavages
\end{itemize}


\section{Connections}

\begin{itemize}
    \item North \& Weingast as a commitment problem -- if you cannot create trust between the people and the state, then the state cannot sufficiently utilize the people as actors/agents of the state
\end{itemize}


\section{My Critiques}

\begin{itemize}
    \item
\end{itemize}


\end{document}
